\chapter{Chapter 4: From Twenty to One Thousand}\label{chapter-4-from-twenty-to-one-thousand}

\emph{In which we discover that the engine is not a toy, that the bottleneck is collision detection, and that cooperation gets stronger at scale.}

\begin{center}\rule{0.5\linewidth}{0.5pt}\end{center}

\section{The Scaling Problem}\label{the-scaling-problem}

Every experiment in Chapters 5 through 11 uses 20-24 agents. This is typical of agent-based modeling --- NetLogo tutorials use 50, Mesa examples use 100, and published ABM papers rarely exceed a few hundred. But a claim about cooperation as ``thermodynamic inevitability'' carries an implicit assumption: that the result holds at scale. Twenty agents in a 100$\times$100 arena is a sparse, simple system. Does cooperation survive density?

\section{The Spatial Hash}\label{the-spatial-hash}

The engine's main computational bottleneck is the O(n²) harmony resonance loop: every pair of agents must be checked for proximity and resonance. At N=20, this is 190 pairs --- trivial. At N=1,000, it is 499,500 pairs per tick at 64 Hz --- 32 million checks per second.

We implemented a grid-based spatial hash (Chapter 1 of the code, \texttt{spatial\_hash.rs}) with cell size equal to the harmony range. This guarantees that all potential cooperation partners are in the same cell or an adjacent cell --- a 3$\times$3 neighborhood in 2D, 3$\times$3$\times$3 in 3D. The check reduces from O(n²) to O(n$\times$k), where k is the average number of agents per cell neighborhood.

Seven unit tests verify correctness, including 3D operation, negative coordinates, and a 1,000-agent population test confirming that queries return a proper subset of the total population.

\section{The Benchmark}\label{the-benchmark}

We benchmarked both brute-force and spatial-hash versions at population sizes from 10 to 1,000:

\begin{longtable}[]{@{}llll@{}}
\toprule\noalign{}
N & Brute Force (ms/tick) & Spatial Hash (ms/tick) & Speedup \\
\midrule\noalign{}
\endhead
\bottomrule\noalign{}
\endlastfoot
10 & 0.02 & 0.04 & 0.5$\times$ (overhead) \\
50 & 0.64 & 0.62 & 1.0$\times$ \\
100 & 2.67 & 2.40 & 1.1$\times$ \\
200 & 14.17 & 7.88 & 1.8$\times$ \\
500 & 165.58 & 59.69 & 2.8$\times$ \\
1,000 & (skipped) & 444.20 & --- \\
\end{longtable}

The spatial hash provides meaningful speedup above N=200 (1.8$\times$) and becomes essential at N=500 (2.8$\times$). At N=1,000, brute force would take approximately 2,500 ms/tick (estimated from the O(n²) trend) --- the spatial hash achieves 444 ms/tick.

However, 444 ms/tick is far from real-time (which would require \textless16 ms at 64 Hz). The bottleneck is not the harmony loop --- it's the physics world step. GJK collision detection on 1,000 spheres dominates the tick time. The spatial hash fixed the harmony bottleneck; the physics broadphase (LBVH) handles collision pairs, but EPA penetration resolution on all overlapping spheres remains expensive.

For the book's purposes, 444 ms/tick at N=1,000 means a full 8,000-tick experiment takes approximately 60 minutes in release mode. This is feasible for research (we ran it) but not for real-time gameplay.

\section{Finding 37: Cooperation at Scale}\label{finding-37-cooperation-at-scale}

The defining experiment ran 100, 250, 500, and 1,000 agents with spatial hashing, 4,000 ticks, five seeds each:

\begin{longtable}[]{@{}lllll@{}}
\toprule\noalign{}
N & Survival & Clusters & Mean Cluster Size & Wall Time \\
\midrule\noalign{}
\endhead
\bottomrule\noalign{}
\endlastfoot
100 & 64\% & 13.6 & 4.8 & 17s \\
250 & 49\% & 13.2 & 9.5 & 103s \\
500 & 74\% & 17.2 & 22.2 & 262s \\
\textbf{1,000} & \textbf{96\%} & \textbf{13.8} & \textbf{72.4} & \textbf{738s} \\
\end{longtable}

Three observations:

\textbf{Cooperation strengthens at scale.} Survival increases from 64\% at N=100 to 96\% at N=1,000. More agents means more potential cooperation partners within harmony range, which means more resonance regeneration per agent. The cooperation benefit scales super-linearly with density.

\textbf{Cluster count is stable.} The number of clusters stays at approximately 14 regardless of population. This is a structural property of the arena: with 4+ energy wells distributed across a 400$\times$400 space, the well geometry creates 3-4 natural basins of attraction, and agents within each basin form 3-4 sub-clusters within harmony range. The cluster count reflects arena geometry, not population.

\textbf{Cluster size scales linearly.} Mean cluster size is approximately N/14 --- each cluster absorbs a proportional share of the population. There is no hard ``Dunbar number'' (a maximum group size beyond which clusters split). The limit is spatial, not cognitive: how many agents fit within harmony range of each other near a well.

The N=250 dip (49\% survival) is the ``scaling valley'': enough agents to create well competition, not enough density for universal cooperation. At N=100, agents are sparse enough that each finds a well without much competition. At N=1,000, agents are dense enough that cooperation partners are always nearby. At N=250, agents compete for wells but can't always find partners --- the worst of both regimes.

\section{Determinism}\label{determinism}

All results in this book are deterministically reproducible. A lock-in test (\texttt{tests/determinism\_lockin.rs}) runs a 500-tick simulation with 8 agents twice using the same seed and asserts bitwise identical final positions, energies, and cooperation counts. Both the lock-in test and a different-seeds-different-results test pass.

This means every experiment can be reproduced exactly:


\begin{verbatim}
\ExtensionTok{cargo} run \AttributeTok{{-}{-}example} cooperation\_1000 \AttributeTok{{-}{-}release}
\end{verbatim}


The same seed produces the same output, every time, on the same platform. (Cross-platform determinism is not guaranteed due to floating-point implementation differences, but single-platform reproducibility is verified.)

\begin{center}\rule{0.5\linewidth}{0.5pt}\end{center}

\emph{Next: Chapter 5 --- Cooperation as Thermodynamic Necessity}
